\chapter{Introducción}
\label{cap:introduccion}

La construcción de una solución para una tarea de procesamiento del lenguaje natural (PLN) requiere elegir datos, modelos, configuraciones y criterios de evaluación. Cada decisión consume tiempo y recursos, y una mejora observada en desarrollo puede desaparecer al cambiar de partición o de dominio. Automatizar la propuesta de alternativas solo es útil si las ejecuciones se verifican, los resultados se conservan y las comparaciones respetan un presupuesto común.

Esta tesis estudia un marco de trabajo basado en agentes para organizar ese proceso. El marco separa tres funciones: proponer candidatos, ejecutarlos y comprobarlos, y analizar la evidencia acumulada para decidir el siguiente paso. Su objeto es la \emph{búsqueda de soluciones}; los agentes no sustituyen al detector o generador que finalmente se evalúa. La contribución se juzga por la calidad de los candidatos encontrados, el costo de encontrarlos, la trazabilidad de las decisiones y la posibilidad de adaptar el proceso a otra tarea.

\section{Caso de estudio y alcance}

El clickbait en noticias ofrece dos tareas relacionadas y distintas. Una publicación puede plantear una pregunta o prometer una revelación sin dar la respuesta; detectar ese mecanismo es una tarea de clasificación. Producir una respuesta breve sustentada en la noticia, o \emph{spoiler}, es una tarea de generación \citep{mordecki2025,ta1c2025}. El corpus TA1C y la tarea compartida de IberLEF 2025 proporcionan datos en español para ambos problemas. La imagen puede incorporarse como entrada opcional en ciertas configuraciones, pero no define el marco de agentes.

La aplicación a detección y \emph{spoiling} permite examinar si una misma interfaz admite objetivos, métricas y restricciones diferentes. La transferencia se estudia en una tarea textual ajena al clickbait, con detección de fraude en mensajes como candidata. La elección final requiere un corpus identificable, etiquetas claras y una partición de evaluación defendible. Una prueba adicional puede aportar evidencia de portabilidad, sin demostrar utilidad para todas las tareas de PLN.

\section{Objetivos y preguntas de investigación}

El objetivo general es diseñar y evaluar un marco de agentes que explore soluciones de PLN de forma reproducible y bajo un presupuesto explícito. Se persiguen tres objetivos específicos: definir una interfaz de tarea reutilizable; comparar la búsqueda asistida por agentes con referencias convencionales en detección y generación; y medir el esfuerzo y el desempeño al adaptar el marco a otra tarea textual.

Las preguntas que organizan la evaluación son:
\begin{enumerate}
  \item ¿Qué calidad alcanzan las soluciones encontradas por el marco frente a búsquedas convencionales con el mismo presupuesto de cómputo y de consultas?
  \item ¿Qué aportan, por separado, la propuesta de candidatos, la ejecución verificada y el análisis de resultados con memoria experimental?
  \item ¿Qué componentes se reutilizan al cambiar de tarea, qué adaptaciones se requieren y cómo cambia el costo total de búsqueda?
\end{enumerate}

\section{Organización de la tesis}

El capítulo~\ref{cap:marco} reúne los antecedentes sobre búsqueda automatizada, agentes y clickbait. El capítulo~\ref{cap:agentes} especifica el marco y su interfaz de tarea. Los capítulos~\ref{cap:metodologia} y~\ref{cap:resultados} documentan el estudio de base realizado con TA1C; sus resultados describen los datos y las líneas de referencia, sin medir todavía el aporte del marco de agentes. El capítulo~\ref{cap:evaluacion-marco} define la comparación y la prueba de transferencia. El capítulo~\ref{cap:conclusiones} sintetiza lo establecido y los límites de la evidencia disponible.
